\def\myname{Рунов К.А.}
\def\mygroup{ИУ7-54Б}
\def\mytheme{Разработка программного обеспечения для моделирования упругих столкновений объектов в пространстве}

\begin{titlepage}
    \newgeometry{left=2cm, right=1cm, top=2.5cm, bottom=2.5cm}
    \fontsize{12pt}{12pt}\selectfont

    \noindent
    \begin{center}
        \begin{minipage}{0.14\textwidth}
            \includegraphics[width=\linewidth]{img/bmstu_logo.jpg}
        \end{minipage}
        \hfill
        \begin{minipage}{0.85\textwidth}\centering\bfseries
            {
                \linespread{1}\selectfont
                \vspace{0.1cm}
                {Министерство науки и высшего образования~Российской~Федерации}

                {Федеральное~государственное~бюджетное~образовательное~учреждение высшего~образования}

                {
                    <<Московский~государственный~технический~университет

                имени~Н.~Э.~Баумана

                (национальный~исследовательский~университет)>>
            }

                {(МГТУ им. Н.~Э.~Баумана)}
                \vspace{0.1cm}
            }
        \end{minipage}

        \vspace{0.2cm}
        \rule{\linewidth}{2.8pt}
        \rule[3ex]{\linewidth}{1pt}

        \begin{flushleft}
            {ФАКУЛЬТЕТ \uline{<<Информатика и системы управления>> \hfill}}

            \vspace{0.5cm}

            {КАФЕДРА \uline{<<Программное обеспечение ЭВМ и информационные технологии>> \hfill}}
        \end{flushleft}

        \vspace{2cm}

        {
            \Large{\textbf{РАСЧЕТНО-ПОЯСНИТЕЛЬНАЯ ЗАПИСКА}}

            \vspace{0.8cm}

            \large{\textbf{\textit{К КУРСОВОЙ РАБОТЕ}}}

            \vspace{0.75cm}

            \large{\textbf{\textit{НА ТЕМУ:}}}
        }

        \vspace{0.5cm}

        % \fontsize{14pt}{14pt}\selectfont

        \setstretch{1.5}

        \begin{tabularx}{\linewidth}{@{}X@{}}
            \uline{
                \large Разработка \large программного \large обеспечения \large для \large моделирования \large упругих \large столкновений \large объектов \large в \large пространстве.
                \hfill}
        \end{tabularx}

        % \begin{tabular}{p{\textwidth}}
        %     \uline
        %     {
        %         {Разработка программного обеспечения для моделирования упругих столкновений объектов в пространстве.}\hfill \,
        %     }
        %     \rule{\linewidth}{0.4pt}
        %     \rule{\linewidth}{0.4pt}
        %     \rule{\linewidth}{0.4pt}
        %     \rule{\linewidth}{0.4pt}
        % \end{tabular}

        % { \large <<\mytheme>> }

        \vfill

        \begin{flushleft}
            {Студент \underset{(Группа)}{\uline{\makebox[3cm][c]{ИУ7-54Б}}}} \hfill
            \underset{(Подпись, дата)}{\uline{\makebox[3cm][c]{}}}
            \underset{\hfill (И. О. Фамилия)}{\uline{\makebox[4cm][r]{К. А. Рунов}}}
            
            \vspace{1.0cm}

            {Руководитель курсовой работы} \hfill
            \underset{(Подпись, дата)}{\uline{\makebox[3cm][c]{}}}
            \underset{\hfill (И. О. Фамилия)}{\uline{\makebox[4cm][r]{А. А. Павельев}}}

            \vspace{1.0cm}
        \end{flushleft}

        \the\year\ г.

    \end{center}
\end{titlepage}

